\chapter{Mental Rotation and Spatial Reasoning}
\label{sec:p19-mental-rotation}

Response time, error rates and eye movements constrain explanations of
a task; none directly displays an internal representation.
This chapter's bounded corpus is Shepard and Metzler's original
three-page shape-identity experiment and Xu and Franconeri's later
feature-binding study \cite{ShepardMetzler1971MentalRotation,
XuFranconeri2015FeatureCapacity}.
The latter is read in its complete PMC author-manuscript version,
NIHMS681047: Experiments 1A, 1B, 1C and 2, including both brief
footnote controls. It changes the task to maintaining features
attached to parts during transformation. These are complementary
constraints on explanation, not two estimates of one common effect.
The reporting standard is \cref{def:p19-evidence-record}.
This component does not complete the visual/spatial reasoning
literature or establish a current consensus.

\section{The geometric task and the observed measure}
\label{sec:p19-rotation-task}

\begin{hbdefinition}{def:p19-rotation-benchmark}{Identity, disparity and response time}
In the original task, \emph{same} means that a rigid rotation can
bring the stipulated three-dimensional objects into congruence.
A reflection is not itself a rotation; the chosen chiral mirror
objects therefore require \emph{different} responses.
This benchmark concerns the objects generating the drawings,
not a claim that every two-dimensional drawing uniquely identifies
a three-dimensional object.
For a stipulated single-axis rotation, disparity \(\delta\) is
measured in degrees. A shorter disparity over one full turn is
\(\min(\theta,360^\circ-\theta)\), \(0\leq\theta\leq360^\circ\).
Measured response time runs between specified stimulus and response
events. It can include encoding, selecting a strategy, transformation,
comparison, decision and motor response. A duration assigned to
one unobserved operation is a model quantity.
\end{hbdefinition}

\section{Shepard and Metzler's 1971 experiment}
\label{sec:p19-shepard-metzler}

The original reports eight adult participants, each receiving
1,600 pairs over eight to ten approximately one-hour sessions.
It does not report a recruitment frame, individual ages, sex
composition, payment or prospective sample-size calculation.
A Stanford affiliation does not establish that the participants
were Stanford undergraduates. Repeated trials do not increase
the independent participant count beyond eight
\cite[pp.~701--703]{ShepardMetzler1971MentalRotation}.

\paragraph{Stimuli and design.}
Ten unfamiliar objects comprised ten cubes joined face to face,
with three right-angle elbows. They formed two five-object sets,
paired across sets by reflection and chosen without nontrivial
within-set rotational or reflection symmetries.
Eighteen computer-generated views per object represented
\(20^\circ\) steps about the vertical axis. Seven views were
selected to avoid complete occlusion of parts and permit two
pairs at each disparity from \(0^\circ\) through \(180^\circ\).
The 70 drawings were reproduced and attached to cards.
The original credits Jih-Jie Chang and Bell Laboratories'
graphics program and recorder; it does not supply every object
coordinate or stimulus card.

There were 400 unique same pairs, each presented twice:
20 depth and 20 picture-plane pairs at each of ten disparities.
Depth pairs used different perspective views. Picture-plane
pairs used differently oriented copies of one view.
Another 400 unique different pairs, also repeated twice,
combined reflection and rotation and were randomly intermixed.
Mirror foils were intended to make local-feature shortcuts less
useful; their selection does not prove that every shortcut was
eliminated. Figure~1 shows an \(80^\circ\) example of each same
type and an incongruent different pair.
Blocks contained at most 200 presentations and were either
pure rotation-type blocks or mixed blocks.

\paragraph{Instructions and timing.}
Participants pulled a right-hand lever for same and a left-hand
lever for different, responding quickly while minimizing errors.
A tone preceded display onset by 0.5 seconds; display onset started
the timer, and the lever stopped it and removed the display.
The drawings' maximum extent was about 4--5 cm, viewing distance
about 60 cm, and center separation \(9^\circ\), through two
apertures on a black surface. The final discussion describes
the pure blocks as making axis and direction predictable.

\paragraph{All reported result types.}
Mean error rate was 3.2 percent, ranging from 0.6 to 5.7 percent
across participants. The displayed time analyses used correct
responses, 96.8 percent overall; errors were said to show a
similar pattern but were not separately tabulated.
Figure~2 plots correct same-response means separately for
picture-plane and depth trials. Both increase approximately
linearly with disparity. Individual means went from about
one second at zero disparity to about four to six seconds at
\(180^\circ\). The bars are described as conservative standard
errors based on the eight participant component means, not
confidence intervals or variability across independent trials.
The article supplies no exact numerical table of those bars.

Polynomial regressions were fitted separately for each participant
and rotation type. All 16 linear components had reported
\(p<0.001\) against deviations from linearity; all quadratic
or higher-order tests had \(p>0.05\). Failure to detect curvature
does not establish an exactly linear population function.
The authors also describe linear patterns across individual objects.
Plane and depth functions had similar slopes and intercepts;
a small apparent depth advantage at large angles was absent
or reversed in four of eight participants and was considered doubtful.

Different pairs had overall mean time 3.8 seconds, nearly a second
longer than same pairs. Since no congruent rotation exists for
those shapes, the paper does not apply the same disparity-to-
congruence function to them. Post-task reports described aligning
one end and discovering that the other ends disagreed.
All participants described imagining rotations, with an explicit
author caution about introspection.
The approximate \(60^\circ\)-per-second interpretation is conditional
on attributing the fitted slope to a rotation process.

Pure-block times remained linear and were no more than 20 percent
lower than mixed-block times. The authors tentatively suggested
that about 80 percent of a typical time might reflect transformation
rather than preparation or search. These were not separately
measured stage durations. Predictability can alter several processes;
the result does not uniquely identify an additive decomposition.
No confidence interval for the rate or decomposition is printed
\cite[p.~703]{ShepardMetzler1971MentalRotation}.

\section{A fully worked fictional calculation}
\label{sec:p19-rotation-worked-example}

\begin{hbtheorem}{thm:p19-rotation-fit-and-identification}{proved=annex:p19.rotation.fit.identification}{Affine arithmetic and a limit of process identification}
For \(m\) stipulated summaries \((\delta_i,t_i)\) with
\(S_{\delta\delta}=\sum_i(\delta_i-\bar\delta)^2>0\), the unique
least-squares affine fit is
\[
 b=\frac{\sum_i(\delta_i-\bar\delta)(t_i-\bar t)}
          {S_{\delta\delta}},\qquad a=\bar t-b\bar\delta.
\]
For milliseconds and degrees, \(b\) has units ms/degree.
If an additive constant-rate model assigns every angle-dependent
mean component to rotation, its rate is \(1000/b\) degrees/second,
provided \(b>0\).
The following summaries are entirely fictional:
\[
\begin{array}{c|rrrrr}
\delta&0&45&90&135&180\\\hline
t\ (\mathrm{ms})&1200&1920&2650&3310&4060\\
\widehat t&1206&1917&2628&3339&4050\\
t-\widehat t&-6&3&22&-29&10
\end{array}
\]
They give \(\bar\delta=90\), \(\bar t=2628\),
\(S_{\delta\delta}=20250\), cross-product sum 319950,
\(b=15.8\), \(a=1206\), and residual sum of squares
\(1470\ {\rm ms}^2\). The model-implied rate is
\(5000/79\simeq63.29\) degrees/second.
These are not recovered study observations and supply no sampling
standard error without further records and assumptions.

On an angle grid \(\delta=j\Delta\), a continuous operation
lasting \(b\delta\) and \(j\) discrete operations each lasting
\(b\Delta\) produce identical conditional time distributions
when given the same intercept and noise. Thus that relation
alone does not identify which internal process occurred.
The construction concerns the stated time observations; it does
not claim that either toy process explains every error or eye movement.
\end{hbtheorem}

\section{Xu and Franconeri's feature-binding task}
\label{sec:p19-xu-franconeri}

The later question concerns retaining color-to-part correspondences,
using simple two-dimensional crosses. Its main outcomes are
change detection, a reported capacity \(K\), and eye movements;
they are not measurements of the 1971 angle--time slope.
The author manuscript reports paid participants aged 18--35,
normal or corrected vision and written consent. Sample sizes
were set in advance by reference to conventional visual-cognition
samples, without a reported power calculation or probability
recruitment frame. Overlap among Experiment~1 pools is unspecified;
Experiment~2 explicitly used a new set
\cite[Experiments~1A--2]{XuFranconeri2015FeatureCapacity}.

\paragraph{Shared stimulus and procedure.}
Experiment~1A used Experiment Builder on a PC, a 17-inch CRT at
75 Hz and \(1024\times768\) pixels, viewed from approximately
56 cm. The display subtended \(32.6^\circ\times24.4^\circ\).
Four cross parts had reported length \(10.7^\circ\) and width
\(2.4^\circ\), with adjacent parts at \(90^\circ\).
Four colors were randomly assigned without replacement:
orange \((233,122,0)\), green \((0,166,0)\), aqua
\((0,136,233)\) and magenta \((230,0,230)\), on grey
\((80,80,80)\). Initial tilt was \(10^\circ\) clockwise or
counterclockwise with equal probability.

A self-initiated trial used a 2400-ms pinwheel cue with mechanical
sound specifying rotation direction and rate. The object then
appeared for 500 ms. During a blank interval of 800, 1600 or
2400 ms, participants imagined it moving at the cued rate and
direction while hearing the sound. The instructions used a
curtain metaphor, although display transitions were immediate.
At test they judged both orientation and color assignments.
Half the images were correct. Rotation foils either swapped
two colors or used an incorrect orientation. In the no-rotation
control, the cue was static and silent and only swaps were foils.
Exact incorrect-orientation offsets are not printed in the
acquired manuscript.

Four distinct consonants preceded each cue for mental rehearsal.
Recall was tested unpredictably on 25 percent of trials;
errors incurred feedback and a three-second penalty.
This was intended to reduce verbal support for the visual task,
not to prove that every verbal strategy was absent.
Guided examples explained the foil types, followed by self-paced
practice and auditory feedback on experimental judgments.

\paragraph{Experiment 1A.}
Twelve participants completed the analyzed sample, including two
replacements for people below 75 percent on the verbal task.
There were 180 randomized trials in five blocks of 36, about
60 minutes. The source crosses direction/task, correct/incorrect
image, six swaps, condition-dependent foil types and three
intervals; the foil factors are conditional on task.
Hit rates were 0.58 for incorrect orientation, 0.64 for swaps
during rotation and 0.81 for swaps without rotation.

The paper computes \(K\) from swap hits and false alarms using
a cited estimator, but does not print its algebra or participant
response tables in the acquired manuscript. The following are
the reported estimates, not a reconstructed fit:
\[
\begin{array}{l|rrl}
\text{Condition}&K&\mathrm{SE}&95\%\ \mathrm{CI}\\\hline
\text{1A no rotation}&1.86&0.15&[1.52,2.19]\\
\text{1A object rotation}&0.96&0.10&[0.74,1.18]\\
\text{1B needle rotation}&0.90&0.16&[0.56,1.25]\\
\text{1B object rotation}&0.97&0.09&[0.76,1.17]\\
\text{1C scaling}&1.63&0.12&[1.35,1.90]
\end{array}
\]
For 1A the capacity comparison was
\(F(1,11)=32.04,p<0.001,\eta^2=0.74\).
Reported accuracy was 76 percent (SE 2.1 percentage points,
CI [71.3,80.7]) during rotation and 87.5 percent
(SE 2.1, CI [82.9,92.1]) without it:
\(F(1,11)=19.63,p=0.001,\eta^2=0.64\).
The interval effect was \(F(2,22)=2.98,p=0.07\).
Printed means for 1600 and 2400 ms were 1.64
(SE 0.13, CI [1.34,1.93]) and 1.21
(SE 0.15, CI [0.88,1.54]), with
\(t(11)=2.29,p=0.04\) for that comparison.
No unreported multiplicity adjustment is added.
Verbal accuracy averaged 90 percent.

\paragraph{Experiment 1B.}
Twelve participants completed 288 randomized trials in eight
blocks of 36, about 90 minutes. In one condition a grey needle
rotated around a static colored cross, with a ticking-sound cue.
In the other the needle was attached to the object's top part
and moved with the object. Both required detecting orientation
errors and swaps, with the same verbal task.
The source crosses two conditions, two directions, correct/
incorrect images, 12 foil levels and three intervals.
This tests the cost of a rotation operation while the remembered
color-bearing object itself remains static.

Orientation hits were 0.65 for object and 0.78 for needle
rotation; swap hits were 0.64 and 0.55.
The capacity comparison was \(F(1,11)=0.54,p=0.48\).
Accuracy was 77.2 percent for needle rotation
(SE 2.4 points, CI [72.0,82.5]) and 77 percent for object
rotation (SE 1.6, CI [73.4,80.6]):
\(F(1,11)=0.02,p=0.89\).
The interval test was \(F(2,22)=0.88,p=0.43\);
verbal accuracy was 90.4 percent. These nonsignificant
comparisons do not establish equivalence.

\paragraph{Experiment 1C and the footnote controls.}
Twelve participants completed 1C; one below the verbal-task
threshold was excluded, leaving eleven analyzed.
The cue showed a noise texture expanding or shrinking by
60 percent over 2400 ms. Participants imagined the cross
scaling at that rate, then judged scale and feature assignments.
Methods specifies a 2400-ms blank interval; Figure~1B's shared
timing label lists the three intervals used elsewhere.
The report gives 144 randomized trials, about 60 minutes,
and factors of direction, image correctness and twelve foil
types, without the complete cell-to-trial repetition schedule.

Scaling hits were 0.50 and swap hits 0.79.
The comparison with 1A rotation detection was
\(t(21)=0.18\), reported nonsignificant.
Scaling \(K=1.63\) exceeded each of the three rotation
capacities above: all \(t(21)>3.56,p<0.003\).
Its comparison with 1A no rotation was \(t(21)=1.15,p=0.26\).
Verbal accuracy was 92.3 percent.
These are cross-experiment comparisons, with different pools
and a specifically fixed interval in the 1C Methods.

Footnote~2 reports an easier-scaling follow-up: scaling hits
0.62, swap hits 0.74, false alarms 0.10 and \(K=1.42\).
The comparison with 1C had \(p=0.32\), and each comparison
with the two object-rotation capacities had \(p=0.02\).
It gives no sample size, interval or full protocol for that
follow-up. Footnote~1 describes a separate swap-only control
with correct rotation angles: false alarms were 0.17,
compared with 0.18 in 1A and 0.16 in 1B.
That control is explicitly not otherwise reported, with no
sample denominator. Both brief reports are retained as such
\cite[Experiments~1A--1C; notes~1--2; Fig.~1]{XuFranconeri2015FeatureCapacity}.

\section{Experiment 2: selection and eye movements}
\label{sec:p19-rotation-eye-movements}

After examining the earlier behavioral data, the authors proposed
that people preferentially selected the topmost part.
Sixteen new participants were eye-tracked with an Eyelink 1000
Remote. The object had two orange/green parts.
Figure~2A gives the 2400-ms cue, 500-ms image, three blank
intervals and test until response. Starting configurations
were top/left, top/right, bottom/left and bottom/right.
The text refers to 1A/1B procedures without a new trial/block
count; neither 180 nor 288 should be invented as that denominator.

Task accuracy, described as an average of the two hit rates,
was 78.8 percent. Eye-analysis exclusions were no eye movements
(8.2 percent), no valid pre-rotation movements (13 percent),
initial saccades outside a roughly \(2.6^\circ\)-radius central
region (0.7 percent), and final saccades during a blink
(3.1 percent) or offscreen (0.04 percent).
The reported retained fraction was 74.9 percent across sixteen
participants. These rounded fractions do not supply exact counts.
The pre-rotation selection analysis further used correct
mental-rotation trials, 80.3 percent of valid trials.

Swap hits were 0.93 and the reported \(K\) was 1.9.
The paper explicitly says this two-part experiment was not
designed to estimate capacity. Every swap involves both parts;
remembering one binding can therefore support perfect swap
detection, as the counting example below proves.

The final pre-rotation fixation/landing location was classified
into four quadrants bounded by diagonals.
Top-quadrant selections were 72.3 and 73.5 percent for
top/left and top/right displays; each comparison with another
quadrant had \(t(15)>3.96,p<0.002\).
For bottom/left and bottom/right displays, left and right
selections were 62.8 and 74.6 percent, respectively;
reported comparisons had \(t(15)>2.61,p<0.02\).

The saccade-position centroid was the arithmetic mean position
during the blank interval. Coordinates were rotated to put
each trial's initial selection at the top before aggregation.
Centroids were below the selection on 71.8 percent of clockwise
and 72.4 percent of counterclockwise trials, both
\(t(15)>4.9,p<0.001\); their direction comparison was
\(t(15)=1.24,p=0.23\).
The horizontal direction comparison was
\(t(15)=10.16,p<0.001\).
Clockwise centroids lay to the right 88.8 percent of the time,
\(t(15)=9.1,p<0.001\); counterclockwise centroids lay to
the left 85.8 percent, \(t(15)=10.58,p<0.001\).

For intended rotations \(40^\circ,80^\circ,120^\circ\),
mean first-to-last fixation angles were
\(27.5^\circ,45.7^\circ,57.4^\circ\).
The Greenhouse--Geisser-corrected repeated-measures test was
\(F(1.3,19.5)=21.05,p<0.001,\eta^2=0.58\);
all reported pairwise tests had \(t(15)>3.5,p<0.003\).
Across participants, eye-angle increment per intended rotation
degree correlated with angular-foil performance:
\(r=0.70,p=0.003\).
Figure~2D labels the latter measure \(d'\); the manuscript
does not give its computation or extreme-rate correction.
An association between tracking and performance does not
identify an effect of experimentally changing eye movements.

The retrospective top/non-top contrasts are printed as:
\[
\begin{array}{l|rr}
&\text{Top involved}&\text{Top not involved}\\\hline
\text{1A hit rate}&82.4\%&62.5\%\\
\mathrm{SE}\text{ (points)}&3.6&5.5\\
95\%\ \mathrm{CI}&[74.6,90.2]&[50.4,74.6]\\
\text{1B hit rate}&78.2\%&40.3\%\\
\mathrm{SE}\text{ (points)}&3.3&6.3\\
95\%\ \mathrm{CI}&[71.0,85.5]&[26.5,54.1]
\end{array}
\]
Their tests were \(F(1,11)=7.02,p=0.02,\eta^2=0.39\)
and \(F(1,11)=93.86,p<0.001,\eta^2=0.90\).
Figure~1D separates task conditions, whereas this paragraph
reports one mean pair per experiment. Those means should not
be relabelled as rotation-only cells.
The authors describe non-top performance as near chance.
No direct test against a specified guessing model is printed,
and the 1A interval begins at 50.4 percent. The numerical
contrast does not establish complete unavailability of every
unselected binding
\cite[Experiment~2; Figs.~1D and 2]{XuFranconeri2015FeatureCapacity}.

\section{A counting benchmark and the comparison's limits}
\label{sec:p19-rotation-capacity-benchmark}

\begin{hbtheorem}{thm:p19-feature-swap-benchmark}{proved=annex:p19.feature.swap.benchmark}{A stipulated detection model}
Suppose \(N\geq2\) distinct-feature parts have exactly
\(k\in\{0,\ldots,N\}\)
perfectly remembered bindings and a uniformly selected
unordered pair is swapped. A swap involving a remembered
part is detected certainly; otherwise a different response
is guessed with probability \(g\), the same as on unchanged trials.
Then
\[
 d_k=1-\frac{\binom{N-k}{2}}{\binom N2},\qquad
 h=d_k+(1-d_k)g,\qquad f=g.
\]
A binomial coefficient with fewer than two available parts
is zero. For \(g<1\), \(d_k=(h-f)/(1-f)\).
For four parts and \(k=0,1,2,3,4\), the detection probabilities
are \(0,1/2,5/6,1,1\). For two parts, remembering either
one or both gives certainty. Thus ceiling performance need
not identify remembering all bindings.
These are explicit ideal assumptions, not a reconstruction
of the published \(K\) estimator.
\end{hbtheorem}

The later authors interpret rotation and needle rotation as
requiring a narrow attentional focus, while scaling permits
attention to the whole object. The eye data add an observation
consistent with following one selected part. The scaling
control and easier version constrain a simple general-difficulty
account; they do not prove all dual-task explanations impossible.
Abstract shape axes, attention marking an axis, and information
supplied by verbal codes or long-term knowledge remain explanatory
proposals. Rotation in depth and educational training suggestions
were not tested in these experiments. The discussion explicitly
notes that a requested \(180^\circ\) turn could be solved by
discrete position swaps, unlike its continuously cued task
\cite[General Discussion]{XuFranconeri2015FeatureCapacity}.

Together the papers distinguish shape-identity timing from the
availability of arbitrary feature bindings. The earlier result
does not establish that every visual detail rotates as one
rigid internal picture. The later result neither refutes the
earlier angle--time relation nor proves a universal one-feature
limit for all spatial representations. Both require task,
population, practice and measurement qualifications.
The geometric answer supplies a normative benchmark;
observed speed and accuracy describe performance.
Faster responding is not thereby universally better reasoning,
and an educational prescription requires additional intervention
and transfer evidence.
